\clearpage
\section{Solutions to the exercises}
\label{sec-solutions-8-applis-representations}

The solutions translate geometric group actions into character values and
then into decomposition or kernel information. The final examples connect
quaternionic Fourier coordinates, cube invariants, and the polynomial
symmetries of self-dual codes.

\begin{correction}{exo-orthogonalite-caracteres}
Row orthogonality says $\Phi K\Phi^*=|G|I$. The square matrix
$U=\Phi K^{1/2}/\sqrt{|G|}$ therefore satisfies $UU^*=I$ and is unitary.
The identity $U^*U=I$ yields
\[
 \Phi^*\Phi=|G|K^{-1},\qquad
 \sum_i\overline{\chi_i(C_j)}\chi_i(C_l)
    =\frac{|G|}{|C_j|}\delta_{jl}.
\]
The conjugation cannot be omitted: characters of a cyclic group of order
three already provide nonreal entries.
\end{correction}

\begin{correction}{exo-repr-groupe-diedral}
Let $r$ rotate the $xy$-plane through $2\pi/n$ and fix the $z$-axis,
and let $s=\operatorname{diag}(1,-1,-1)$. These matrices are proper
rotations satisfying $r^n=s^2=I$ and $srs=r^{-1}$. They give the
specified $2n$ rotations. Their characters are
\[
 \chi(r^k)=1+2\cos(2\pi k/n),\qquad \chi(sr^k)=-1.
\]
The complexified $xy$-plane carries $\rho_1$, while the $z$-axis
carries $\psi_2$, which is $1$ on rotations $r^k$ and $-1$ on
$sr^k$. Thus the decomposition is $\rho_1\oplus\psi_2$ for every
$n\geq3$, whether $n$ is odd or even. Both multiplicities are one.
\end{correction}

\begin{correction}{exo-representation-s3}
The cycle $(123)$ rotates through $2\pi/3$. The transposition $(12)$
reflects about the line at angle $\pi/3$, whereas $(23)$ fixes the real
axis. Consequently,
\[
 \rho((123))=\frac12\begin{pmatrix}-1&-\sqrt3\\\sqrt3&-1\end{pmatrix},\quad
 \rho((12))=\frac12\begin{pmatrix}-1&\sqrt3\\\sqrt3&1\end{pmatrix},\quad
 \rho((23))=\begin{pmatrix}1&0\\0&-1\end{pmatrix}.
\]
In the class order identity, transposition, three-cycle, the class sizes
are $1,3,2$ and the character table is
\[
 \begin{array}{c|rrr}
 &1&(12)&(123)\\\hline
 1&1&1&1\\\varepsilon&1&-1&1\\\chi_\rho&2&0&-1
 \end{array}.
\]
The last row has squared norm $(4+2)/6=1$, proving irreducibility.
The degree squares sum to $1+1+4=6$, proving completeness.
\end{correction}

\begin{correction}{exo-action-face-cube}
In the class order $1,(12),(123),(1234),(12)(34)$, the vertex and
edge permutation characters are respectively
\[
 \chi_V=(8,0,2,0,0),\qquad \chi_E=(12,2,0,0,0).
\]
A nontrivial rotation fixes vertices only when its axis is a body
diagonal, in which case it fixes its two endpoints. A half-turn about
an axis joining opposite edge midpoints fixes those two edges as
unordered pairs, although it exchanges their endpoints. No other
nonidentity rotation fixes an edge. Taking weighted inner products
with the character table gives
\[
 V\simeq 1\oplus\varepsilon\oplus s\oplus W,\qquad
 E\simeq 1\oplus s^{\oplus2}\oplus W\oplus W'.
\]
The dimensions are $1+1+3+3=8$ and $1+6+3+2=12$.
\end{correction}

\begin{correction}{exo-caracteres-s4}
A double transposition $g$ has order two, so its representing matrix
is diagonalizable with eigenvalues in $\{1,-1\}$. Trace two in
dimension two forces both eigenvalues to be one. The representation
is therefore trivial on
$H=\{1,(12)(34),(13)(24),(14)(23)\}$.
A representation trivial on $H$ defines a well-defined representation
of $G/H$ by $gH\mapsto\rho(g)$; pullback reverses this construction
and preserves invariant subspaces.

The action on $12|34$, $13|24$, and $14|23$ has kernel $H$.
Its image contains a transposition and a three-cycle, and hence is
all of $\mathfrak S_3$. Thus $\mathfrak S_4/H\simeq\mathfrak S_3$.
The pullback of the standard two-dimensional representation has
character $(2,0,-1,0,2)$ and is $W'$.
\end{correction}

\begin{correction}{exo-representation-gpe-simple}
\begin{enumerate}
\item The first $p-1$ entries of a tuple in $X$ determine the last,
so $|X|=|G|^{p-1}$. Cyclic rotation preserves $X$: the product after
one rotation is $g_1^{-1}(g_1\cdots g_p)g_1=1$. Each orbit has size
one or $p$. The fixed tuples are $(g,\ldots,g)$ with $g^p=1$.
Their number is therefore congruent to $|X|$ modulo $p$. It is a
positive multiple of $p$, since the identity gives one fixed tuple
and $p\mid|X|$. There is a nonidentity $g$ with $g^p=1$, necessarily
of order $p$.
\item If $\rho$ had degree two, divisibility of irreducible degrees
would imply $2\mid|G|$, and the preceding argument would give an
element $t$ of order two.
\item The kernel is normal, so simplicity and nontriviality of
$\rho$ make it trivial. The determinant homomorphism also has a
normal kernel. If this kernel were trivial, $G$ would embed in the
abelian group $\CC^*$ and be abelian. Hence $\det\rho(g)=1$ for all
$g$. The minimal polynomial of $\rho(t)$ divides $X^2-1$, so this
matrix is diagonalizable with eigenvalues $\pm1$. Determinant one
and faithfulness exclude $I$ and force $\rho(t)=-I$. Since this
matrix commutes with every $\rho(g)$, faithfulness makes $t$ central.
The center of a nonabelian simple group is trivial, a contradiction.
\end{enumerate}
\end{correction}

\begin{correction}{exo-grpe-quaternionique}
\begin{enumerate}
\item The matrices $L(i),L(j),L(k)$ satisfy the displayed multiplication
table. Bilinearity then gives $L(qq')=L(q)L(q')$ for all $q,q'$.
The first column of $L(q)$ recovers $q$, so this map is injective;
associativity of matrices proves associativity of quaternion multiplication.
The map $\RR[H_8]\to\mathbb H$ sends the formal basis vector $\delta_g$
to the quaternion $g$. Its kernel has basis
\[
 \delta_{-g}+\delta_g,\qquad g=1,i,j,k.
\]
It is the ideal generated by $\delta_{-1}+\delta_1$ and has dimension
four. In the group algebra, $\delta_{-1}$ is an independent basis
vector; it becomes $-\delta_1$ only in this quotient.
\item Direct multiplication verifies the same relations for $M$ and gives
\[
 M(q)^*M(q)=(a^2+b^2+c^2+d^2)I.
\]
Thus the eight group elements have unitary images. A common invariant
complex line would be an eigenspace of $M(i)=\operatorname{diag}(\imath,-\imath)$,
hence a coordinate line, but $M(j)$ exchanges these two lines. The
representation is irreducible, with character $(2,-2,0,0,0)$ on
the classes $\{1\},\{-1\},\{\pm i\},\{\pm j\},\{\pm k\}$.
\item Every one-dimensional character factors through
$H_8/\{\pm1\}\simeq\FF_2^2$. The signs assigned to $i$ and $j$
are independent and their product is the sign assigned to $k$.
These four characters and $M$ are complete because $4\cdot1^2+2^2=8$.
In the prescribed element order, the rows of the following matrix $A$
are an orthonormal basis for the normalized inner product on $\CC[H_8]$:
\[
 A=\operatorname{diag}(1,1,1,1,\sqrt2,\sqrt2,\sqrt2,\sqrt2)
 \begin{pmatrix}
 1&1&1&1&1&1&1&1\\
 1&1&1&1&-1&-1&-1&-1\\
 1&1&-1&-1&1&1&-1&-1\\
 1&1&-1&-1&-1&-1&1&1\\
 1&-1&\imath&-\imath&0&0&0&0\\
 0&0&0&0&-1&1&-\imath&\imath\\
 0&0&0&0&1&-1&-\imath&\imath\\
 1&-1&-\imath&\imath&0&0&0&0
 \end{pmatrix}.
\]
The last four rows are $\sqrt2$ times the four matrix entries of $M$.
In particular, $AA^*=8I$.
\item Products of these basis functions give an orthonormal basis on
$H_8^n$, which has $8^n$ points. After fixing a compatible product
ordering, the normalized Fourier coefficients of a signal $f$ are
\[
 c=8^{-n}\,\overline A^{\otimes n}f.
\]
For Euclidean unitary coordinates, use instead
$U^{\otimes n}f$ with $U=\overline A/\sqrt8$.
The conjugation comes from taking inner products with the basis
functions. Successive transforms along the $n$ array axes compute
the tensor product, as in Program~\ref{tensor-product-listing}.
\end{enumerate}
\end{correction}

\begin{correction}{exo-ring-invariants}
A cube rotation is a signed coordinate permutation with determinant
one. The sum of coordinate squares is therefore invariant. If its
underlying permutation is $\sigma$ and its three signs are
$\epsilon_1,\epsilon_2,\epsilon_3$, then
$\epsilon_1\epsilon_2\epsilon_3=\operatorname{sgn}(\sigma)$.
Consequently $XYZ$ transforms by $\operatorname{sgn}(\sigma)$ and
its square is invariant.

Expanding the four factors gives
\[
 F=2(X^2Y^2+X^2Z^2+Y^2Z^2)-X^4-Y^4-Z^4.
\]
This expression is invariant under all signed coordinate permutations,
so both $F$ and $F^2$ are invariant. The polynomial $Q$ is alternating
in $X^2,Y^2,Z^2$; it also transforms by $\operatorname{sgn}(\sigma)$.
Thus $Q^2$ and $XYZ\,Q$ are invariant under proper cube rotations.
The four planes occurring in $F$ are perpendicular to the four body
diagonal lines. Factoring $Q$ into $X\pm Y$, $X\pm Z$, and $Y\pm Z$
gives six planes perpendicular to face-diagonal directions.
\end{correction}

\begin{correction}{exo-codes-autoduaux}
\begin{enumerate}
\item Since $\dim\Cc^\perp=n-k$, self-duality gives $n=2k$.
\item MacWilliams gives
$W_\Cc(X,Y)=2^{-k}W_\Cc(X+Y,X-Y)$. Homogeneity of degree $n=2k$
absorbs the factor as $2^{-n/2}$, proving invariance under $A$.
\item Put $t=\sqrt2-1$, $\ell=X+tY$, and $m=Y-tX$.
The substitution $A$ fixes $\ell$ and negates $m$. Every invariant
is therefore a polynomial in $\ell,m^2$, which are algebraically
independent. The quadratic $q=Y(X-Y)$ is invariant and is not a
multiple of $\ell^2$. Since the invariant quadratics have basis
$\ell^2,m^2$, it follows that $m^2$ is a linear combination of
$\ell^2,q$ with a nonzero coefficient of $q$. Thus
$\CC[X,Y]^{G_1}=\CC[\ell,q]$, again with independent generators.
\item For $c\in\Cc=\Cc^\perp$, the identity
$c\cdot c=\sum_i c_i=0$ in $\FF_2$ makes its weight even.
Hence $W_\Cc(X,-Y)=W_\Cc(X,Y)$. The matrix $AD$ is rotation
through $\pi/4$, of order eight, and $D$ is a reflection conjugating
it to its inverse. The resulting group has eight rotations and
eight reflections. Although $-I=(AD)^4$ belongs to this group,
invariance under $-I$ alone does not impose even powers of $Y$.
\item In the coordinates $z=X+\imath Y$, $w=X-\imath Y$, the
rotation multiplies $z,w$ by inverse primitive eighth roots and
$D$ exchanges them. A rotation-invariant monomial has exponent
difference divisible by eight. After imposing the exchange symmetry,
all invariants are generated by $t=zw$ and $u=z^8+w^8$:
the sums $z^{8r}+w^{8r}$ satisfy the recurrence with initial values
$2,u$ and characteristic polynomial $T^2-uT+t^8$.
Direct expansion gives
\[
 t=P=X^2+Y^2,\qquad u=2P^4-64Q,
 \quad Q=X^2Y^2(X^2-Y^2)^2.
\]
These generators are algebraically independent. For example, the
Jacobian of $(zw,z^8+w^8)$ is $8(w^8-z^8)$, a nonzero polynomial.
Thus $\CC[X,Y]^{G_2}=\CC[P,Q]$, in agreement with the Molien series
$1/((1-t^2)(1-t^8))$.
Every homogeneous self-dual weight enumerator consequently has the form
\[
 W_\Cc=\sum_{j=0}^{\lfloor n/8\rfloor}
       c_j P^{n/2-4j}Q^j.
\]
For the extended Hamming code, expansion gives
$P^4-4Q=X^8+14X^4Y^4+Y^8$, the enumerator found earlier.
\end{enumerate}

Program~\ref{listing-invariant-polynomials} implements Reynolds averages
of all monomials of degrees one through the group order. The helper
\texttt{exponents\_of\_degree} enumerates their exponent tuples, and
\texttt{matrix\_group} in the same module constructs the finite group
by exact matrix multiplication. In characteristic zero, Noether's
bound makes the output a generating set. It is generally redundant:
normalizing and removing duplicate polynomials does not remove all
algebraic dependencies. The two generators $P,Q$ above are obtained
by the additional mathematical argument, not by an automatic minimality
claim for this routine.
\begin{lstlisting}[language=Python]
import sympy as sp
from invariants import matrix_group, reynolds_generators
X, Y = sp.symbols("X Y")
A = sp.Matrix([[1, 1], [1, -1]]) / sp.sqrt(2)
D = sp.diag(1, -1)
generators = reynolds_generators(matrix_group([A, D]), (X, Y))
\end{lstlisting}
\begin{listing}[H]
\PythonListing{invariants.py}{21}{27}
\PythonListing{invariants.py}{30}{41}

\caption{Exact Reynolds averages giving an overcomplete set of invariant generators.}
\label{listing-invariant-polynomials}
\end{listing}
\end{correction}
